\documentclass[12pt]{article}
\usepackage[utf8]{inputenc}
\usepackage[T1]{fontenc}
\usepackage[margin=1in]{geometry}
\usepackage{setspace}
\usepackage{xcolor}
\usepackage[colorlinks=true, urlcolor=blue!60!black, linkcolor=blue!60!black]{hyperref}
\setlength{\parindent}{2em}
\setlength{\parskip}{1.25em}
\renewcommand{\baselinestretch}{1.0}
\begin{document}
\begin{center}
{\LARGE \textit{West Virginia v.\ EPA}, 597 U.S.\ 697 (2022)\par}
\vspace{0.8em}
{\large Case Brief --- Four-Page Edition\par}
\end{center}
\vspace{1em}

\section{Memory Jogger}

Naming the ``major questions doctrine'' in a majority opinion for the first time, the Court held 6--3 that \S~111(d) of the Clean Air Act did not clearly authorize EPA's Clean Power Plan, which set carbon-dioxide limits premised on shifting electricity generation away from coal (\href{https://supreme.justia.com/cases/federal/us/597/20-1530/}{Justia}).

\section{Detailed Case Facts}

Section 111(d) is a ``gap-filler'' requiring EPA to set guidelines for existing stationary sources based on the ``best system of emission reduction'' (BSER) that is ``adequately demonstrated,'' with States enforcing the resulting limits. Before 2015 EPA had used it only a handful of times, always basing BSER on technology that makes a plant run more cleanly. The 2015 Clean Power Plan (CPP) set BSER through three ``building blocks'': heat-rate improvements at coal plants, ``generation shifting'' from coal to gas, and shifting to wind and solar. Operators could comply by reducing coal generation, investing in cleaner plants, or buying credits in a cap-and-trade regime. EPA projected coal's share of electricity would fall from 38\% in 2014 to 27\% by 2030, and no coal plant could comply through on-site technology alone.

\section{Procedural History}

Twenty-seven States sued, and the Supreme Court took the unprecedented step of staying the CPP in 2016. The Trump-era EPA repealed it in 2019 and adopted the narrower Affordable Clean Energy (ACE) Rule. The D.C.\ Circuit vacated both the repeal and ACE, rejecting the major questions doctrine. After the new administration said it would write a new rule, West Virginia, North Dakota, and coal companies obtained certiorari.

\section{Judicial Votes}

\textbf{6--3}, reversed and remanded (June 30, 2022).
\begin{itemize}
\item \textbf{Majority:} Chief Justice Roberts, joined by Justices Thomas, Alito, Gorsuch, Kavanaugh, and Barrett.
\item \textbf{Concurring:} Justice Gorsuch, joined by Justice Alito.
\item \textbf{Dissenting:} Justice Kagan, joined by Justices Breyer and Sotomayor.
\end{itemize}

\section{Holding}

The case was justiciable: the D.C.\ Circuit's vacatur purported to revive the CPP, injuring the States, and EPA's voluntary cessation did not moot the dispute. On the merits, \S~111(d) does not authorize emissions caps based on generation shifting. In ``extraordinary cases,'' where the ``history and the breadth of the authority'' asserted and its ``economic and political significance'' give reason to hesitate, a merely plausible textual basis is not enough; the agency must point to ``clear congressional authorization.''

\section{Analysis of Opinions}

\textbf{Majority (Roberts).} Synthesizing \textit{Brown \& Williamson}, \textit{Gonzales v.\ Oregon}, \textit{UARG}, \textit{Alabama Realtors}, and \textit{NFIB v.\ OSHA}, the Court grounded the doctrine in separation of powers and legislative intent. Four signs triggered it: EPA found a ``transformative expansion'' of power in a statutory ``backwater''; generation shifting worked a ``fundamental revision of the statute''; grid reliability and energy prices lie outside EPA's expertise; and Congress had ``considered and rejected'' cap-and-trade. The word ``system'' could not supply clear authorization, and the Acid Rain and NAAQS trading programs operate under caps Congress or health standards set. The Court did not decide whether \S~111(d) is limited to inside-the-fenceline measures.

\textbf{Gorsuch (concurring).} The doctrine is a substantive clear-statement rule protecting Article~I's vesting of legislative power. Triggers include political significance, a significant share of the economy, and intrusion on state domains; clues include statutory structure, age, past interpretations, and expertise mismatch.

\textbf{Kagan (dissenting).} The Court issued an advisory opinion on a defunct rule and abandoned textualism. ``Best system of emission reduction'' is broad, not vague; Congress omitted ``technological'' for existing sources; and cap-and-trade is an ordinary pollution-control system within EPA's wheelhouse.

\textbf{Later developments.} \textit{Biden v.\ Nebraska} (2023) applied the doctrine to student-loan cancellation, with Justice Barrett recasting it as contextual textualism. In \textit{Learning Resources v.\ Trump} (2026), the Court held that IEEPA did not authorize President Trump's tariffs, splitting four ways over the doctrine.

\section{Examples: Future Applications}

\begin{enumerate}
\item \textbf{Same-Side: AI Licensing Under an Old Statute.} An agency uses a 1974 lab-safety statute to license frontier AI after Congress failed to act---an unheralded, costly expansion lacking clear authorization.
\item \textbf{Same-Side: Ratemaking to Retire Coal.} Declaring all coal-plant rates after 2035 unjust would restructure the electricity mix through a general provision.
\item \textbf{Opposite-Side: Express AI Statute.} A new statute expressly authorizing AI licensing supplies clear authorization; any challenge shifts to nondelegation.
\item \textbf{Opposite-Side: Source-Specific Rule.} Section 111(d) guidelines based on on-site technology claim no unheralded power.
\item \textbf{Fence-Sitter: Climate Disclosure.} SEC climate-risk disclosure rules sit within its expertise but carry large costs on a charged subject; Barrett's and Gorsuch's approaches may diverge.
\end{enumerate}

\section{Critique}

\textbf{Canon or context?} The Gorsuch--Barrett--Kagan debate asks whether demanding clear authorization when a broad term covers the action is faithful reading or a judge-made avoidance rule. \textbf{Partisan ratchet?} Critics called the doctrine deregulatory, but \textit{Learning Resources} saw Roberts, Gorsuch, and Barrett apply it against a Republican President. \textbf{Failed bills.} Congressional inaction may reflect gridlock or belief that the agency already had authority.

\section{Key Quotations}

\begin{quote}
``Thus, in certain extraordinary cases, both separation of powers principles and a practical understanding of legislative intent make us `reluctant to read into ambiguous statutory text' the delegation claimed to be lurking there. \dots{} To convince us otherwise, something more than a merely plausible textual basis for the agency action is necessary. The agency instead must point to `clear congressional authorization' for the power it claims.''
\end{quote}
\noindent --- Chief Justice Roberts, for the Court, 597 U.S.\ at 723.

\begin{quote}
``But of course almost anything could constitute such a `system'; shorn of all context, the word is an empty vessel. Such a vague statutory grant is not close to the sort of clear authorization required by our precedents.''
\end{quote}
\noindent --- Chief Justice Roberts, for the Court.

\begin{quote}
``The current Court is textualist only when being so suits it. When that method would frustrate broader goals, special canons like the `major questions doctrine' magically appear as get-out-of-text-free cards.''
\end{quote}
\noindent --- Justice Kagan, dissenting.

\end{document}
